%% ============================================================
\section{Introduction}
\label{sec:intro}

The C programming language remains dominant in systems software,
embedded systems, and security-critical infrastructure.  Despite decades
of research in programming language safety, C's lack of memory safety
guarantees continues to produce vulnerabilities: buffer overflows, null
pointer dereferences, integer overflows, and use-after-free defects
account for the majority of CVEs in C codebases~\cite{cwe-top25}.

The SEI CERT C Coding Standard~\cite{cert-c} codifies 120+ rules and
200+ recommendations spanning 17 categories, from integer handling to
concurrency to POSIX-specific interfaces.  While the standard is widely
referenced in secure coding guidelines (e.g., MISRA C, IEC 62443),
tool support for comprehensive CERT C checking remains fragmented.
Commercial tools such as Coverity~\cite{coverity}, Klocwork, and
Polyspace implement subsets of CERT C rules alongside proprietary
checkers, but their closed-source nature limits reproducibility.
Open-source alternatives like cppcheck~\cite{cppcheck} and
clang-tidy~\cite{clang-tidy} implement approximately 20 CERT C checks
each, leaving the majority of the standard uncovered.

This paper presents SqC (\textbf{S}oftware Code \textbf{Q}uality for
\textbf{C}), a static analysis tool that implements 311 CERT C rules
across all 17 categories.  The paper describes the tool's architecture,
the iterative process used to reduce false positives, and the results of
evaluating SqC against both a standardized benchmark and real-world
codebases.

The rest of this paper is organized as follows.  \Cref{sec:background}
covers the CERT C standard and the Juliet benchmark.
\Cref{sec:architecture} describes SqC's architecture.
\Cref{sec:methodology} details the testing methodology.
\Cref{sec:results} presents benchmark results.
\Cref{sec:fp-reduction} walks through the iterative false positive
reduction process.  \Cref{sec:worked-example} provides a worked
example.  \Cref{sec:limitations} discusses limitations.
\Cref{sec:related} covers related work.  \Cref{sec:conclusion}
concludes.
